\section*{Chapter \ref{ch:b2:huckel} --- H\"uckel Theory and Conjugated Systems}

\begin{solution}{exo:b2:huckel:1}
Allyl levels ($n = 3$): $\alpha + 1.414\beta$, $\alpha$, $\alpha - 1.414\beta$. Cation (2 electrons):
$2\alpha + 2.828\beta$; radical (3): $3\alpha + 2.828\beta$; anion (4): $4\alpha + 2.828\beta$. The extra
electrons go into the nonbonding level and change only the $\alpha$ part.
\end{solution}

\begin{solution}{exo:b2:huckel:2}
$2(\alpha + 1.618\beta) + 2(\alpha + 0.618\beta) = 4\alpha + 4.472\beta$.
\end{solution}

\begin{solution}{exo:b2:huckel:3}
Benzene (6), cyclopentadienyl anion (6), tropylium (6): \omterm{def:b2:huckel:aromatic}{aromatic}. Cyclopentadienyl
cation (4), cyclobutadiene (4), planar cyclo-octatetraene (8): \omterm{def:b2:huckel:aromatic}{antiaromatic} if planar
(real cyclo-octatetraene is tub-shaped and \omterm{def:b2:huckel:aromatic}{non-aromatic}).
\end{solution}

\begin{solution}{exo:b2:huckel:4}
An octagon in the circle: levels $\alpha + 2\beta$, $\alpha + 1.414\beta$ (twice), $\alpha$ (twice),
$\alpha - 1.414\beta$ (twice), $\alpha - 2\beta$. Eight electrons: two in the lowest, four in the
first pair, one in each orbital of the nonbonding pair.
\end{solution}

\begin{solution}{exo:b2:huckel:5}
$4.472\beta - 4\beta = 0.472\beta$, that is $0.472 \times 75 = \qty{35}{kJ/mol}$, against about
\qty{10}{kJ/mol} from the hydrogenation data: H\"uckel's model, with $\beta$ fitted to benzene,
overestimates the conjugation of open chains.
\end{solution}

\begin{solution}{exo:b2:huckel:6}
$p_{12} = p_{34} = 0.894$, $p_{23} = 0.447$: the end bonds C1--C2 and C3--C4 are the
shortest, the central bond longer but shorter than a single bond.
\end{solution}

\begin{solution}{exo:b2:huckel:7}
Anion, six electrons: $2 \times 2 + 4 \times 0.618$, $E_\pi = 6\alpha + 6.472\beta$, a closed shell:
\omterm{def:b2:huckel:aromatic}{aromatic}. Cation, four: $2 \times 2 + 2 \times 0.618$, $4\alpha + 5.236\beta$, with one electron
in each orbital of the degenerate pair: \omterm{def:b2:huckel:aromatic}{antiaromatic}, very unstable.
\end{solution}

\begin{solution}{exo:b2:huckel:8}
Seven-membered ring: $\alpha + 2\beta$, $\alpha + 1.247\beta$ (twice), then antibonding levels. Six
electrons, as in the cation \ce{C7H7+}, fill the bonding levels exactly: an \omterm{def:b2:huckel:aromatic}{aromatic}
cation. The compound ionises readily into tropylium and bromide.
\end{solution}

\begin{solution}{exo:b2:huckel:9}
$1.802 + 1.247 + 0.445 - 0.445 - 1.247 - 1.802 = 0$: the sum of the energies is $6\alpha$.
\end{solution}

\begin{solution}{exo:b2:huckel:10}
$n = 6$: $m_k = 2\cos(k\pi/7) = 1.802$, $1.247$, $0.445$, $-0.445$, $-1.247$, $-1.802$. Six
electrons: $E_\pi = 6\alpha + 2(1.802 + 1.247 + 0.445)\beta = 6\alpha + 6.988\beta$;
delocalisation $0.988\beta$.
\end{solution}

\begin{solution}{exo:b2:huckel:11}
In the square, the two electrons of highest energy sit one in each orbital of a
degenerate nonbonding pair. Stretching two opposite bonds lifts the degeneracy:
one orbital goes down, the other up, and both electrons pair in the lower one, which
lowers the energy. The molecule settles as a rectangle with two localised double
bonds; it remains \omterm{def:b2:huckel:aromatic}{antiaromatic} and extremely reactive.
\end{solution}

\begin{solution}{exo:b2:huckel:12}
Matrix (in units of $\beta$, with $\alpha$ as origin) $\begin{pmatrix}0 & 1\\ 1 & 1\end{pmatrix}$:
$m = 1.618$ and $-0.618$, $E = \alpha + 1.618\beta$ (bonding) and $\alpha - 0.618\beta$. For the
\omterm{def:b2:diatomic-mos:bonding}{bonding orbital} $c_X = 1.618c_C$, so $c_C^2 = 0.276$, $c_X^2 = 0.724$; with two electrons,
$q_C = 0.55$, $q_X = 1.45$: the $\pi$ density is drawn to X.
\end{solution}

\begin{solution}{pb:b2:huckel:1}
\textbf{1.} $-123.1 - (-4.3) = \qty{-118.8}{kJ/mol}$.
\textbf{2.} $-123.1 - 104.6 = \qty{-227.7}{kJ/mol}$.
\textbf{3.} $-123.1 - 82.9 = \qty{-206.0}{kJ/mol}$.
\textbf{4.} $3 \times (-118.8) = \qty{-356.4}{kJ/mol}$.
\textbf{5.} $356.4 - 206.0 = \qty{150}{kJ/mol}$.
\textbf{6.} The diene releases $2 \times 118.8 - 227.7 = \qty{10}{kJ/mol}$ less than two
isolated double bonds: conjugation of an open chain gives a small gain, the closed
ring of benzene fifteen times more.
\textbf{7.} A $6 \times 6$ matrix with 1 between each atom and its two neighbours (1--2, 2--3,
\dots, 6--1) and 0 elsewhere.
\textbf{8.} $m = 2\cos(2\pi k/6)$: $2$, $1$, $1$, $-1$, $-1$, $-2$: $E = \alpha + 2\beta$, $\alpha + \beta$ (twice),
$\alpha - \beta$ (twice), $\alpha - 2\beta$.
\textbf{9.} Two electrons in $\alpha + 2\beta$, four in the pair $\alpha + \beta$.
\textbf{10.} $E_\pi = 6\alpha + 8\beta$.
\textbf{11.} $3(2\alpha + 2\beta) = 6\alpha + 6\beta$.
\textbf{12.} $2\beta$, that is $2|\beta|$ of stabilisation.
\textbf{13.} $2 + 1 + 1 - 1 - 1 - 2 = 0$: the six levels add up to $6\alpha$.
\textbf{14.} $2|\beta| = \qty{150}{kJ/mol}$: $|\beta| = \qty{75}{kJ/mol}$.
\textbf{15.} $0.472 \times 75 = \qty{35}{kJ/mol}$, against \qty{10}{kJ/mol}: the model fitted
on benzene overestimates the open chain.
\textbf{16.} The occupied orbitals can be taken as $\psi_1 = (1, 1, 1, 1, 1, 1)/\sqrt6$,
$\psi_2 = (2, 1, -1, -2, -1, 1)/\sqrt{12}$ and $\psi_3 = (0, 1, 1, 0, -1, -1)/2$. For the bond
1--2: $p_{12} = 2(\frac16 + \frac{2}{12} + 0) = \frac23$; by symmetry every bond has $p = 2/3$.
\textbf{17.} Between the end (0.894) and central (0.447) bonds of butadiene.
\textbf{18.} The six bonds are equivalent by symmetry; their length, \qty{139.7}{pm}, lies
between the double bond of ethene (\qty{133.9}{pm}) and the single bond of ethane
(\qty{153.6}{pm}).
\textbf{19.} $0.988|\beta| = \qty{74}{kJ/mol}$, half of benzene's: closing the ring doubles the
\omterm{def:b2:huckel:delocalisation}{delocalisation energy}.
\textbf{20.} $\alpha + 2\beta$, $\alpha + 1.414\beta$ (twice), $\alpha$ (twice), $\alpha - 1.414\beta$ (twice),
$\alpha - 2\beta$.
\textbf{21.} The last two electrons go one in each orbital of the nonbonding pair: an open
shell, \omterm{def:b2:huckel:aromatic}{antiaromatic}.
\textbf{22.} It folds into a tub: the p orbitals of neighbouring double bonds are no longer
parallel, the bonds alternate, and the molecule behaves as a polyene.
\textbf{23.} Four electrons: the same half-filled pair, worse because the molecule cannot
avoid planarity; it distorts into a rectangle and is extremely reactive.
\textbf{24.} A planar monocyclic conjugated system is \omterm{def:b2:huckel:aromatic}{aromatic} with $4N + 2$ $\pi$ electrons,
\omterm{def:b2:huckel:aromatic}{antiaromatic} with $4N$.
\textbf{25.} $\boldsymbol{|\beta| \approx \qty{75}{kJ/mol}}$.
\end{solution}
